% !TEX root = ../main.tex
\chapter{圏：対象・射・合成・恒等射}
\label{chap:ch11_category}

\begin{chapterabstract}
本章から圏論に入る。最初に扱うのは、圏という最も基本的な概念である。
圏とは、対象と射があり、射を合成でき、何もしない射がある世界である。
本章では、一般圏論のすべてを一度に形式化するのではなく、まず「型と関数の世界」を使って、対象、射、恒等射、合成、恒等律、結合律を理解する。
Lean 4 では、恒等関数と関数合成を小さく定義し、パイプラインの括弧の付け方を変えても結果が変わらないことを確認する。
\end{chapterabstract}

\begin{learninggoals}
\item 圏を「合成できる処理の世界」として説明できる。
\item 対象、射、恒等射、合成を、型、関数、何もしない処理、パイプラインとして読み替えられる。
\item 恒等律と結合律を、リファクタリングやパイプライン設計の安全性として理解できる。
\item Lean 4 で恒等関数、関数合成、合成の法則を小さな定理として書ける。
\end{learninggoals}

\section{圏論を一言でいうと「合成できるものの理論」}

ソフトウェアでは、処理を一つだけで完結させることは少ない。
入力を検証し、正規化し、変換し、保存し、表示用に整える。
API gateway では request を受け取り、認証し、ログを付け、内部サービスへ流す。
ETL では extract、transform、load を順に並べる。
これらはすべて「処理をつなぐ」作業である。

圏論は、この「つなげる」という行為を中心に置く。
重要なのは、個々の処理の中身を細かく見ることではない。
どの処理とどの処理をつなげられるか、つないだ結果が何になるか、何もしない処理を挟んでも意味が変わらないか、三つ以上をつなぐときに括弧の付け方で意味が変わらないかを見る。

\begin{intuitionbox}
圏を最初に読むときは、「合成できる矢印のネットワーク」と考えるとよい。
点があり、点から点へ向かう矢印があり、矢印を順につなぐことができる。
ソフトウェアの言葉では、点は型や状態の形、矢印は関数や変換、矢印をつなぐことはパイプラインである。
\end{intuitionbox}

\FigurePlaceholder{fig-ch11-category-overview.png}{0.90}{対象・射・恒等射・合成をソフトウェアの処理として見る}{fig:ch11-category-overview}

本章では、最も身近な例として「型と関数の世界」を使う。
たとえば、\texttt{String}、\texttt{Nat}、\texttt{Bool} は対象の例である。
\texttt{String} から \texttt{Nat} への関数、\texttt{Nat} から \texttt{String} への関数は射の例である。
そして、二つの関数の入出力型が合えば、それらを合成できる。

\section{ソフトウェア上の困りごと：処理を分けると意味が壊れるのか}

次のような処理を考える。
文字列で来たユーザーIDを解釈し、ユーザー名を取り出し、挨拶文を作る。
実際のシステムでは DB や外部 API が絡むかもしれないが、ここでは単純な関数として見る。

\begin{softwarebox}
処理を分割すると、レビューでは次の問いが出る。
「処理を三つに分けたが、元の意味は保たれているのか」
「先に前半をまとめるか、後半をまとめるかで結果は変わらないのか」
「何もしない adapter や middleware を挟んでも本当に影響はないのか」
圏の法則は、これらの問いを整理するための最小限の語彙になる。
\end{softwarebox}

たとえば、次の三つの処理を考える。

\begin{itemize}
\item \texttt{parseUserId : String -> Nat}
\item \texttt{fetchUserName : Nat -> String}
\item \texttt{renderGreeting : String -> String}
\end{itemize}

これらは、型が合っているので順に合成できる。
\texttt{parseUserId} の出力は \texttt{Nat} であり、\texttt{fetchUserName} の入力も \texttt{Nat} である。
同様に、\texttt{fetchUserName} の出力は \texttt{String} であり、\texttt{renderGreeting} の入力も \texttt{String} である。

ここで、前半を先にまとめるか、後半を先にまとめるかを考える。

\[
  (\texttt{renderGreeting} \circ \texttt{fetchUserName}) \circ \texttt{parseUserId}
\]

\[
  \texttt{renderGreeting} \circ (\texttt{fetchUserName} \circ \texttt{parseUserId})
\]

この二つは、実行順序としてはどちらも「IDを解釈し、名前を取り出し、挨拶文を作る」である。
括弧は「どこを先にひとまとまりの部品として見るか」を表しているだけで、データが流れる順序を変えていない。
これが結合律のソフトウェア的な意味である。

\section{対象を型、射を関数として見る}

圏には、対象と射がある。
対象は点のようなものであり、射は対象から対象へ向かう矢印である。
ただし、圏論で重要なのは、対象そのものの中身をのぞき込むことではない。
対象どうしを結ぶ射と、その射をどう合成するかである。

\begin{definitionbox}
圏の基本要素は次の通りである。
\begin{itemize}
\item \textbf{対象}：矢印の始点や終点になるもの。
\item \textbf{射}：ある対象から別の対象へ向かう矢印。
\item \textbf{恒等射}：各対象にある、何もしない射。
\item \textbf{合成}：つなげられる二つの射から、新しい射を作る操作。
\end{itemize}
本章では、対象を型、射を関数として読む。
\end{definitionbox}

型と関数の世界では、対象は \texttt{String} や \texttt{Nat} のような型である。
射は、\texttt{String -> Nat} のような関数である。
関数 \(f : A \to B\) は、対象 \(A\) から対象 \(B\) への射として読める。
このとき、\(A\) を始域、\(B\) を終域と呼ぶことがある。

\begin{table}[htbp]
\centering
\begin{tabular}{lll}
\toprule
圏論の語彙 & 型と関数での読み方 & ソフトウェアでの例 \\
\midrule
対象 & 型 & \texttt{String}, \texttt{Nat}, \texttt{UserDto} \\
射 & 関数 & parser, formatter, mapper \\
恒等射 & 恒等関数 & 何もしない middleware \\
合成 & 関数合成 & パイプライン、adapter chain \\
恒等律 & no-op を挟んでも同じ & 不要な wrapper の除去 \\
結合律 & 括弧を変えても同じ & パイプラインの段階的グループ化 \\
\bottomrule
\end{tabular}
\caption{第11章で使う圏論語彙の読み替え}
\label{tab:ch11-vocabulary}
\end{table}

表\ref{tab:ch11-vocabulary}で重要なのは、圏論の言葉を暗記することではない。
自分が普段扱っている型、関数、変換、パイプラインを、同じ構造として見られるようになることである。
この見方は、第12章の同型、第14章の関手、第15章の自然変換へ続く。

\section{恒等射は何もしない処理}

各対象には、恒等射がある。
対象 \(A\) の恒等射は、\(A\) から \(A\) へ戻る射であり、何も変えない。
型と関数の世界では、これは恒等関数である。

\[
  \mathrm{id}_A : A \to A
\]

任意の値 \(x : A\) に対して、\(\mathrm{id}_A(x) = x\) である。
つまり、入力をそのまま返す。

\begin{softwarebox}
恒等射は、ソフトウェアでは「何もしない処理」と読める。
たとえば、デバッグ用に一時的に挟んだ middleware、将来の拡張用に置いた adapter、あるいは設定で無効化された変換が、入力をそのまま返すなら恒等射に相当する。
恒等律は、そのような no-op を前後に挟んでも全体の意味が変わらないことを表す。
\end{softwarebox}

ここで注意したいのは、恒等射は「何も書かない」ことではない。
明示的に「入力をそのまま返す」という射である。
プログラム上では、\texttt{id} や \texttt{fun x => x} に相当する。

\section{合成はパイプライン}

射 \(f : A \to B\) と \(g : B \to C\) があるとする。
\(f\) の終域と \(g\) の始域が同じなので、\(f\) の後に \(g\) を実行できる。
この合成を \(g \circ f : A \to C\) と書く。

\[
  (g \circ f)(x) = g(f(x))
\]

右側の式を見ると、\(f\) が先に実行され、その結果に \(g\) が適用される。
表記では \(g \circ f\) と \(g\) が左に来るが、処理順序は \(f\) から \(g\) である。
これは最初につまずきやすい点である。

\begin{warningbox}
\(g \circ f\) は、「先に \(f\)、次に \(g\)」である。
日本語で左から右にパイプラインを読む感覚とは逆に見えることがある。
本書では、式としては標準的な \(g \circ f\) を使い、説明では常に「先に \(f\)、次に \(g\)」と読む。
\end{warningbox}

\FigurePlaceholder{fig-ch11-composition-order.png}{0.90}{\(g \circ f\) を「先に \(f\)、次に \(g\)」として読む}{fig:ch11-composition-order}

合成を使うと、小さな部品を大きな部品として扱える。
\(f : A \to B\) と \(g : B \to C\) を合成した \(g \circ f : A \to C\) は、もう一つの射である。
この「合成したものもまた同じ種類の部品になる」という性質が、圏論をソフトウェア設計に結びつける。
パイプラインの一部をまとめて名前を付け、さらに別の処理と合成できるからである。

\section{圏の二つの法則}

圏には、恒等射と合成があるだけでは不十分である。
それらが自然に振る舞うための法則が必要になる。
本章で使う法則は二つである。
恒等律と結合律である。

\subsection{恒等律}

射 \(f : A \to B\) に対して、\(A\) 側の恒等射と \(B\) 側の恒等射を合成しても、\(f\) は変わらない。

\[
  f \circ \mathrm{id}_A = f
\]

\[
  \mathrm{id}_B \circ f = f
\]

ソフトウェア的には、入力側に no-op を挟んでも、出力側に no-op を挟んでも、全体の振る舞いは変わらないということである。
これは、不要な adapter や wrapper を取り除くリファクタリングの基礎になる。

\subsection{結合律}

射 \(f : A \to B\)、\(g : B \to C\)、\(h : C \to D\) があるとする。
このとき、三つの射を合成する方法には、括弧の付け方として二通りある。

\[
  (h \circ g) \circ f
\]

\[
  h \circ (g \circ f)
\]

結合律は、この二つが等しいことを主張する。

\[
  (h \circ g) \circ f = h \circ (g \circ f)
\]

実行順序は、どちらも \(f\)、\(g\)、\(h\) の順である。
違うのは、どこを先に一つの部品としてまとめるかだけである。

\begin{intuitionbox}
結合律は、「処理をどの粒度で箱詰めしても、外から見た結果は同じ」という法則である。
三つの処理を、一つずつ並べても、前半二つを先にまとめても、後半二つを先にまとめても、入力から出力への全体の変換は同じである。
\end{intuitionbox}

\FigurePlaceholder{fig-ch11-associativity.png}{0.90}{三つの処理の括弧の付け方を変えても実行順序は変わらない}{fig:ch11-associativity}

ここで重要なのは、結合律は「処理の順序を入れ替えてよい」とは言っていない点である。
\(f\)、\(g\)、\(h\) の順序を保ったまま、括弧だけを変えてよいと言っている。
順序を入れ替えてよいかどうかは、別の性質であり、一般には成り立たない。

\section{Lean 4 で関数合成の法則を確認する}

Lean では、型と関数をそのまま扱える。
本章では Mathlib の高度な圏論ライブラリは使わない。
自分で恒等関数と合成を小さく定義し、何が起きているかを見えるようにする。

\begin{leanbox}
ここでの Lean コードは、一般の圏を定義するものではない。
まず、型を対象、関数を射とする特別な世界で、恒等射と合成がどのように見えるかを確認する。
一般の圏や Mathlib の \texttt{CategoryTheory} API は、発展章で扱う。
\end{leanbox}

\begin{listing}[htbp]
\begin{minted}{lean4}
namespace Ch05Category

-- 対象を型、射を関数として見る。
-- idFn は「何もしない処理」である。
def idFn {A : Type} (x : A) : A :=
  x

-- comp g f は「先に f、次に g」を実行する合成である。
def comp {A B C : Type} (g : B -> C) (f : A -> B) : A -> C :=
  fun x => g (f x)

end Ch05Category
\end{minted}
\caption{恒等関数と関数合成を小さく定義する}
\label{lst:ch11_id_comp}
\end{listing}

\texttt{comp g f} は、数学の \(g \circ f\) に対応する。
入力 \texttt{x} に対して、まず \texttt{f x} を計算し、その結果に \texttt{g} を適用する。
したがって、\texttt{comp g f x} は \texttt{g (f x)} と同じになる。

次に、恒等律と結合律を「任意の入力 \texttt{x} に対して結果が同じ」という形で書く。
この段階では、関数そのものの等しさを深追いせず、各入力での出力一致として確認する。
関数そのものの等しさは、第21章の関数外延性で改めて扱う。

\begin{listing}[htbp]
\begin{minted}{lean4}
namespace Ch05Category

-- 左から恒等処理を合成しても、結果は変わらない。
theorem id_left_pointwise {A B : Type} (f : A -> B) (x : A) :
    comp idFn f x = f x := by
  rfl

-- 右から恒等処理を合成しても、結果は変わらない。
theorem id_right_pointwise {A B : Type} (f : A -> B) (x : A) :
    comp f idFn x = f x := by
  rfl

-- 三つの関数を合成するとき、括弧の付け方は結果に影響しない。
theorem comp_assoc_pointwise {A B C D : Type}
    (h : C -> D) (g : B -> C) (f : A -> B) (x : A) :
    comp (comp h g) f x = comp h (comp g f) x := by
  rfl

end Ch05Category
\end{minted}
\caption{恒等律と結合律を点ごとの等式として確認する}
\label{lst:ch11_laws}
\end{listing}

ここで使っている \texttt{rfl} は、両辺を定義に従って展開すると同じ形になることを Lean に確認させる道具である。
たとえば、\texttt{comp idFn f x} を展開すると、\texttt{idFn (f x)} になり、さらに \texttt{f x} になる。
したがって、証明は \texttt{rfl} で終わる。

\section{ミニケース：パイプラインの括弧を変えても意味は変わらない}

ここから、ソフトウェアに近い形の小さなパイプラインを Lean で書く。
次の例では、本物の DB や API は使わない。
証明したいのは外部システムの挙動ではなく、関数合成としてのパイプライン構造だからである。

\begin{listing}[htbp]
\begin{minted}{lean4}
namespace Ch05Category

-- 例を小さくするため、文字列の長さをユーザーIDの代わりに使う。
def parseUserId (s : String) : Nat :=
  s.length

def fetchUserName (id : Nat) : String :=
  "user-" ++ toString id

def renderGreeting (name : String) : String :=
  "hello, " ++ name

-- 前半を後でまとめる見方。
def pipelineLeft : String -> String :=
  comp (comp renderGreeting fetchUserName) parseUserId

-- 後半を先にまとめる見方。
def pipelineRight : String -> String :=
  comp renderGreeting (comp fetchUserName parseUserId)

#eval pipelineLeft "abc"  -- => "hello, user-3"

end Ch05Category
\end{minted}
\caption{三段階の処理パイプラインを定義する}
\label{lst:ch11_pipeline}
\end{listing}

\texttt{pipelineLeft} と \texttt{pipelineRight} は、括弧の付け方が違うだけである。
どちらも、先に \texttt{parseUserId} を実行し、次に \texttt{fetchUserName} を実行し、最後に \texttt{renderGreeting} を実行する。
そのため、任意の入力文字列 \texttt{s} に対して結果は一致する。

\begin{listing}[htbp]
\begin{minted}{lean4}
namespace Ch05Category

-- 点ごとに見ると、両方のパイプラインは同じ式へ展開される。
theorem pipeline_assoc_pointwise (s : String) :
    pipelineLeft s = pipelineRight s := by
  rfl

-- calc を使うと、レビュー用の変形ログとして読める。
theorem pipeline_review_log (s : String) :
    pipelineLeft s = pipelineRight s := by
  calc
    pipelineLeft s
        = renderGreeting (fetchUserName (parseUserId s)) := by rfl
    _   = pipelineRight s := by rfl

end Ch05Category
\end{minted}
\caption{パイプラインの結合律をレビュー可能な形で確認する}
\label{lst:ch11_pipeline_proof}
\end{listing}

\texttt{pipeline\_review\_log} は、数学的には短い証明である。
しかし実務上は、「変更前後のパイプラインは、どちらも同じ中間式に展開される」というレビュー記録として読める。
このような小さな証明は、リファクタリングや設計変更の説明に使いやすい。

\section{恒等処理の除去を証明する}

恒等律は、no-op を取り除くときに使える。
たとえば、何もしない文字列変換 \texttt{noOpString} をパイプラインの末尾に挟んだとする。
この処理は入力をそのまま返すので、取り除いても結果は変わらない。

\begin{listing}[htbp]
\begin{minted}{lean4}
namespace Ch05Category

-- 何もしない文字列処理。
def noOpString : String -> String :=
  idFn

-- no-op を最後に挟んだパイプライン。
def pipelineWithNoOp : String -> String :=
  comp noOpString pipelineLeft

-- 任意の入力で、no-op ありと no-op なしの結果は一致する。
theorem remove_noop_pointwise (s : String) :
    pipelineWithNoOp s = pipelineLeft s := by
  rfl

end Ch05Category
\end{minted}
\caption{何もしない処理を取り除いても結果が変わらない}
\label{lst:ch11_noop}
\end{listing}

この例は小さいが、設計上の意味は大きい。
処理の入口や出口に、将来の拡張用 adapter を置くことがある。
その adapter が本当に何もしないなら、全体の意味は変わらない。
Lean では、その主張を「すべての入力で結果が一致する」という形で書ける。

\FigurePlaceholder{fig-ch11-noop-removal.png}{0.88}{恒等処理をパイプラインから取り除いても結果が変わらない}{fig:ch11-noop-removal}

\section{圏を定義として見る}

ここまで、型と関数の世界で対象、射、合成、恒等射を見てきた。
ここで、圏の定義を少しだけ一般化して眺める。

\begin{definitionbox}
圏とは、次のデータと法則を持つ構造である。
\begin{itemize}
\item 対象の集まり。
\item 各対象 \(A, B\) に対する、\(A\) から \(B\) への射の集まり。
\item 各対象 \(A\) に対する恒等射 \(\mathrm{id}_A : A \to A\)。
\item 射 \(f : A \to B\) と \(g : B \to C\) から射 \(g \circ f : A \to C\) を作る合成。
\item 恒等律と結合律。
\end{itemize}
\end{definitionbox}

この定義は、最初は抽象的に見える。
しかし、型と関数の世界を思い出せば、各項目は自然に読める。
型が対象であり、関数が射であり、恒等関数が恒等射であり、関数合成が合成である。
恒等律と結合律は、関数合成について成り立つ当たり前の法則である。

本章では、一般の圏を Lean の構造体として完全に定義するところまでは進まない。
理由は二つある。
第一に、初学者にとっては、依存型や universe、型クラスが同時に出てくると負担が大きい。
第二に、本書の目的は、まず実務上の変換と合成を読むことである。
したがって、当面は「型と関数の圏」を手で触りながら理解する。

\begin{warningbox}
「圏」は「型と関数」だけを意味するわけではない。
順序、状態遷移、データベーススキーマ、仕様の強弱、証明の変換など、さまざまなものが圏として見られる。
ただし、本章では最初の足場として、型と関数に限定する。
限定することで、対象・射・合成・恒等射を具体的に確認できる。
\end{warningbox}

\section{実務への読み替え}

圏の法則は、実務上の設計判断に直接対応する。
本章では、特に三つの読み替えを押さえる。

\subsection{処理の境界を型で見る}

射を合成できるかどうかは、終域と始域が合うかどうかで決まる。
ソフトウェアでは、これは出力型と入力型が合うかどうかに対応する。
型が合っていなければ、adapter が必要になる。
adapter を入れるなら、その adapter が何を保存するのかを明確にする必要がある。

\subsection{no-op を設計上の部品として扱う}

何もしない処理は、無意味とは限らない。
ログをまだ入れていない middleware、将来の互換層、設定で無効化されたフィルタなど、構造上の場所を確保するために no-op が使われることがある。
恒等律は、それらが本当に no-op である限り、前後に合成しても意味が変わらないことを保証する。

\subsection{括弧の変更と順序変更を区別する}

結合律は、括弧の変更を許す。
しかし、順序変更を許すわけではない。
\texttt{validate -> normalize -> save} と \texttt{normalize -> validate -> save} は、一般には同じではない。
一方で、\texttt{(validate -> normalize) -> save} と \texttt{validate -> (normalize -> save)} は、処理順序が同じなら同じ全体処理として扱える。

\begin{softwarebox}
設計レビューでは、次のように言える。
「この変更は処理順序を入れ替えているのではなく、同じ順序の処理を別の粒度でまとめ直しているだけです。したがって、合成の結合律として見ると意味は変わりません。」
この言い方ができると、リファクタリングの安全性を説明しやすくなる。
\end{softwarebox}

\section{よくある誤解}

\subsection{誤解1：圏論は難しい対象の中身を調べる理論である}

圏論は、対象の中身だけを見る理論ではない。
むしろ、対象どうしの関係、つまり射と合成を見る。
型の内部表現をすべて知ることよりも、その型からどこへ変換できるか、どの変換を合成できるかを重視する。

\subsection{誤解2：結合律があるなら処理を自由に並べ替えられる}

結合律は、順序を変えてよいとは言っていない。
\((h \circ g) \circ f\) と \(h \circ (g \circ f)\) は同じだが、\(g \circ f\) と \(f \circ g\) は型が合わないことすら多い。
型が合ったとしても、結果が同じとは限らない。

\subsection{誤解3：恒等射は実務では役に立たない}

恒等射は、何もしないので軽く見られやすい。
しかし、何もしない処理が「本当に何もしない」と言えることは重要である。
互換層やデフォルト設定、無効化されたフィルタを安全に取り除けるかどうかは、恒等律として考えられる。

\section{次章への接続}

本章では、型を対象、関数を射、関数合成を合成、恒等関数を恒等射として見た。
この見方により、処理パイプラインの括弧を変えても意味が変わらないことや、no-op を挟んでも結果が変わらないことを Lean で確認した。

次章では、二つの型やデータ構造が「違う形をしているが同じ情報を持つ」と言える状況を扱う。
これが同型である。
本章で学んだ恒等射と合成は、同型の round-trip law を読むための基礎になる。
たとえば、\texttt{migrate} してから \texttt{rollback} すると恒等処理になる、という主張は、本章の言葉で言えば「合成が恒等射になる」という主張である。

圏は、対象、射、恒等射、合成を持ち、恒等律と結合律を満たす構造である。

型と関数の世界では、対象は型、射は関数、恒等射は恒等関数、合成は関数合成として読める。

恒等律は、何もしない処理を前後に挟んでも意味が変わらないことを表す。

結合律は、三つ以上の処理を合成するとき、括弧の付け方を変えても全体の意味が変わらないことを表す。

Lean では、これらの性質を小さな関数と定理として確認できる。
